%%% BEGIN fig:disall-seed
\begin{figure}[htbp]
\centering
\begin{tikzpicture}[fd,
  fb lab/.style={font=\footnotesize,inner sep=1.5pt},
  fb op/.style={font=\large},
  fb hregion/.style={hard blob,fill=gray!22,draw=gray!60!black,line width=0.6pt,dash pattern=on 2.5pt off 1.5pt}]
  % ======================= left: one K gluon attached to the hard part =======================
  \begin{scope}
    \coordinate (V) at (-0.4,1.55);       % photon vertex
    \coordinate (G) at (1.0,2.3);         % attachment of the K gluon on the struck quark
    \coordinate (E) at (2.05,2.86);       % outgoing struck quark
    \begin{scope}[on background layer]
      \draw[fb hregion,rotate around={28.2:(0.3,1.92)}] (0.3,1.92) ellipse[x radius=1.12,y radius=0.42];
    \end{scope}
    \node[jet blob,ellipse,minimum width=1.9cm,minimum height=0.7cm,inner sep=0pt] (A) at (0,0) {};
    \node[font=\small] at (-0.45,0) {$A$};
    \draw[photon] (-1.75,2.3) -- (V);
    \draw[quark] (-0.4,0.34) -- (V);
    \draw[quark] (V) -- (G);
    \draw[quark] (G) -- (E);
    \draw[longit] (0.6,0.33) to[out=72,in=-108] (G);
    \fill[blue!60!black] ($(G)+(-0.06,-0.19)$) circle[radius=1.7pt];
    \node[vertex] at (V) {}; \node[vertex] at (G) {};
    \path[mom'={$q$},mom sep=0.24cm,mom length=0.45cm] (-1.75,2.3) -- (V);
    \path[mom={$l$},mom sep=0.24cm,mom length=0.45cm] (-0.4,0.34) -- (V);
    \path[mom={$p'-k$},mom sep=0.3cm,mom length=0.42cm,mom pos=0.42] (V) -- (G);
    \path[mom={$p'$},mom sep=0.24cm,mom length=0.42cm] (G) -- (E);
    \path[mom'={$k$},mom sep=0.26cm,mom length=0.45cm,mom pos=0.45] (0.6,0.33) to[out=72,in=-108] (G);
    \node[fb lab,text=gray!35!black] at ($(G)+(0.42,-0.36)$) {$H$};
    \node[fb lab,anchor=north] at (0.08,-0.42) {$k_\mu H^\mu_a=g\,t^a\,H_0$};
    \node[font=\large,anchor=east] at (-1.75,1.2) {$\displaystyle\sum_{H\text{ 内的附着}}$};
  \end{scope}
  \node[fb op] at (2.75,1.4) {$=$};
  % ======================= middle: H_0 =======================
  \begin{scope}[xshift=4.55cm]
    \coordinate (V) at (-0.25,1.65);
    \draw[photon] (-1.3,2.35) -- (V);
    \draw[quark] (-0.25,0.6) -- (V);
    \draw[quark] (V) -- (1.15,2.4);
    \node[vertex] at (V) {};
    \path[mom={$l+k$},mom sep=0.24cm,mom length=0.45cm] (-0.25,0.6) -- (V);
    \path[mom={$p'$},mom sep=0.24cm,mom length=0.42cm] (V) -- (1.15,2.4);
    \node[fb lab] at (0.6,1.42) {$H_0$};
  \end{scope}
  \node[fb op] at (6.15,1.4) {$\times$};
  % ======================= right: the Wilson-line vertex of the PDF operator =======================
  \begin{scope}[xshift=7.75cm]
    \coordinate (O) at (-0.35,1.65);     % operator vertex psi(0)
    \coordinate (W) at (1.75,2.8);       % towards infinity along n_B
    \coordinate (Y) at ($(O)!0.5!(W)$);
    \node[jet blob,ellipse,minimum width=1.9cm,minimum height=0.7cm,inner sep=0pt] (A) at (0,0) {};
    \node[font=\small] at (-0.45,0) {$A$};
    \draw[quark] (-0.35,0.34) -- (O);
    \draw[wilson=0.78] (O) -- (W);
    \node[opvertex] at (O) {};
    \draw[longit] (0.55,0.33) to[out=70,in=-110] (Y);
    \node[vertex] at (Y) {};
    \path[mom'={$k$},mom sep=0.26cm,mom length=0.45cm,mom pos=0.45] (0.55,0.33) to[out=70,in=-110] (Y);
    \node[fb lab,anchor=east] at ($(O)+(-0.12,0.05)$) {$\psi(0)$};
    \node[fb lab,anchor=west] at (W) {$\infty n_B$};
    \node[fb lab,anchor=north west] at ($(Y)+(0.12,-0.06)$) {$\ii g\,t^a n_B^\alpha$};
    \node[fb lab,anchor=south east] at ($($(O)!0.5!(Y)$)+(-0.02,0.12)$) {$\dfrac{-\ii}{n_B\cdot k-\ii0}$};
    \node[fb lab,anchor=north] at (0.4,-0.42) {$W^{(1)}[\infty n_B,0]$：未来指向};
  \end{scope}
  % summary formula
  \node[font=\footnotesize,anchor=north] at (4.4,-1.15) {$H_\mu\,\dfrac{k^\mu n_B^\alpha}{n_B\cdot k-\ii0}=\Big[\ii g\,t^an_B^\alpha\,\dfrac{-\ii}{n_B\cdot k-\ii0}\Big]H_0=W^{(1)}[\infty n_B,0]\,H_0$};
\end{tikzpicture}
\caption{DIS 中 PDF Wilson 线的种子（VF-WL-8，\cref{prop:method-ward-tree}）：一个来自靶共线子图 $A$ 的 K 胶子（蓝色，流入硬部分的动量为 $k$，圆点标出带标量极化 $k_\mu$ 的硬端）对硬部分 $H$ 的所有附着之和（领头阶即附着在末态被击夸克 $p'$ 上），
经 Ward 恒等式 $k_\mu H^\mu_a=g\,t^aH_0$ 化为去掉该胶子、入射动量为 $l+k$ 的硬部分 $H_0$，乘以夸克场 $\psi(0)$ 处出发、沿 $n_B$ 的未来指向 Wilson 线 $W^{(1)}[\infty n_B,0]$ 的单胶子项（双线箭头为路径方向 $0\to\infty n_B$）。
程函传播子中的 $-\ii0$ 继承自被替代的被击夸克传播子：$(p'-k)^2+\ii0\simeq-2p'^-k^++\ii0$ 的极点位于 $k^+$ 上半平面；Drell--Yan 中同样的附着落在入射反夸克上，给出 $+\ii0$。}
\label{fig:disall-seed}
\end{figure}
%%% END fig:disall-seed
